\documentclass[11pt]{article}
\usepackage[a4paper,margin=1in]{geometry}
\usepackage{amsmath,amssymb,amsthm,mathtools}
\usepackage{graphicx}
\usepackage[hidelinks]{hyperref}
\usepackage{microtype}

\newtheorem{theorem}{Theorem}
\newtheorem{lemma}[theorem]{Lemma}
\newtheorem{corollary}[theorem]{Corollary}
\theoremstyle{remark}
\newtheorem{remark}[theorem]{Remark}

\newcommand{\C}{\mathbb C}
\newcommand{\ns}{\operatorname{ns}}
\newcommand{\supp}{\operatorname{supp}}
\newcommand{\1}{\mathbf 1}

\title{Exact low-dimensional numerical-sparsity uncertainty constants\\
and a stronger asymptotic lower bound}
\author{Candidate partial solution packet for arXiv:1504.01014, Problem 3}
\date{11 August 2026}

\begin{document}
\maketitle

\begin{center}
\fbox{\parbox{0.9\textwidth}{\textbf{Status.}
Candidate substantial partial solution, likely valid.  We determine
$c(1),c(2),c(3)$ exactly and improve the source's asymptotic lower bound by a
factor of $16.2$.  The exact problem for $n\ge4$ remains open.}}
\end{center}

\section{Source problem}

For $x\in\C^n\setminus\{0\}$, Bandeira, Lewis, and Mixon define the
\emph{numerical sparsity}
\[
 \ns(x)=\frac{\|x\|_1^2}{\|x\|_2^2}.
\]
Problem 3 on page 5 of \cite{BLM} asks for the exact value of
\begin{equation}\label{eq:def-c}
 c(n)=\frac1n\sup_{U\in\mathrm U(n)}
       \inf_{x\ne0}\bigl(\ns(x)+\ns(Ux)\bigr).
\end{equation}
The source proves $\liminf c(n)\ge1/540000$ and records the upper bound
$c(n)<1+1/n$ for $n\ge4$.

\begin{figure}[ht]
 \centering
 \includegraphics[width=\textwidth]{figures/open_problem_crop.png}
 \caption{Problem 3 and the original bounds on page 5 of \cite{BLM}.}
\end{figure}

\section{Proof intuition}

In dimensions two and three, the trivial basis-vector upper bound can be
attained by the Fourier matrix.  For a unit vector $z$, numerical sparsity is
one plus twice the sum of the magnitudes of all pairwise coordinate products.
For $F_2$, the sole output pair product is
$(x_0^2-x_1^2)/2$.  For $F_3$, three phase-weighted sums of the three output
pair products isolate, up to unit phases, the quadrics
$x_0^2-x_1x_2$, $x_1^2-x_0x_2$, and $x_2^2-x_0x_1$.  Choosing the largest
coordinate then turns the desired estimate into the elementary inequality
$a(a+b+c)\ge a^2+b^2+c^2$.

For the asymptotic bound, we keep the source's random-unitary restricted
isometry theorem but replace its coarse width lemma.  Sorting a null vector
into decreasing blocks of size $k$ controls the head by RIP and the tail by
the $\ell_1$ norm.  Pythagoras then gives a width constant valid for every
$\delta<1$, allowing a much larger RIP distortion and a much larger $k$.

\section{The first three constants}

Let $F_n$ denote the normalized Fourier matrix
\[
 (F_nx)_j=\frac1{\sqrt n}\sum_{r=0}^{n-1}\omega_n^{jr}x_r,
 \qquad \omega_n=e^{-2\pi i/n}.
\]

\begin{theorem}[Exact low-dimensional values]\label{thm:small}
For the constants in \eqref{eq:def-c},
\[
 \boxed{c(1)=2,\qquad c(2)=\frac32,\qquad c(3)=\frac43.}
\]
The matrices $F_2$ and $F_3$ attain the two nontrivial suprema.
\end{theorem}

\begin{proof}
For every unit vector $v\in\C^n$, Cauchy--Schwarz gives $\ns(v)\le n$.
Testing an arbitrary $U$ at a coordinate vector therefore gives
\begin{equation}\label{eq:upper}
 \inf_{x\ne0}\bigl(\ns(x)+\ns(Ux)\bigr)\le n+1.
\end{equation}
This proves the required upper bounds $3$ and $4$ in dimensions two and
three.  The case $n=1$ is immediate.

For $n=2$, normalize $\|x\|_2=1$ and put $y=F_2x$.  Then
\[
 \ns(x)+\ns(y)
 =2+2\bigl(|x_0x_1|+|y_0y_1|\bigr),
 \qquad y_0y_1=\frac{x_0^2-x_1^2}{2}.
\]
Write $a=\max(|x_0|,|x_1|)$ and $b=\min(|x_0|,|x_1|)$.  The reverse triangle
inequality gives
\begin{align*}
 |x_0x_1|+|y_0y_1|
 &\ge ab+\frac{a^2-b^2}{2}\\
 &\ge \frac{a^2+b^2}{2}=\frac12,
\end{align*}
where the second inequality is $ab\ge b^2$.  Thus the uncertainty sum is at
least $3$, and equality holds at a coordinate vector.  Hence $c(2)=3/2$.

For $n=3$, define
\[
 S(z)=|z_0z_1|+|z_1z_2|+|z_2z_0|.
\]
For unit $z$ one has $\ns(z)=1+2S(z)$.  Put $y=F_3x$ and
$\omega=\omega_3$.  Direct expansion, using $1+\omega+\omega^2=0$, gives
\begin{align}
 y_0y_1+y_1y_2+y_2y_0
   &=x_0^2-x_1x_2,\label{eq:q0}\\
 y_0y_1+\omega y_1y_2+\omega^2y_2y_0
   &=\omega(x_1^2-x_0x_2),\label{eq:q1}\\
 y_0y_1+\omega^2y_1y_2+\omega y_2y_0
   &=\omega^2(x_2^2-x_0x_1).\label{eq:q2}
\end{align}
Each left side has modulus at most $S(y)$.  Relabel the coordinates so that
their magnitudes are $a\ge b\ge c$.  The identity corresponding to the
largest coordinate and the reverse triangle inequality yield
\[
 S(y)\ge a^2-bc.
\]
Consequently,
\[
 S(x)+S(y)
 \ge ab+ac+bc+a^2-bc
 =a(a+b+c)
 \ge a^2+b^2+c^2=1.
\]
It follows that
\[
 \ns(x)+\ns(F_3x)=2+2(S(x)+S(F_3x))\ge4.
\]
Again a coordinate vector gives equality.  Combining with
\eqref{eq:upper} proves $c(3)=4/3$.
\end{proof}

\section{A direct RIP-to-width estimate}

We use the standard definition that a matrix $\Phi$ satisfies
$(k,\delta)$-RIP if
\[
 (1-\delta)\|w\|_2^2\le\|\Phi w\|_2^2
 \le(1+\delta)\|w\|_2^2
\]
for every $k$-sparse $w$.

\begin{lemma}[Nullspace width for every $\delta<1$]\label{lem:width}
If $\Phi$ satisfies $(k,\delta)$-RIP with $0<\delta<1$, then every
$z\in\ker\Phi$ satisfies
\begin{equation}\label{eq:width}
 \|z\|_2^2\le\frac{2}{(1-\delta)k}\|z\|_1^2.
\end{equation}
\end{lemma}

\begin{proof}
Partition the coordinates of $z$ into blocks $S_0,S_1,\ldots$ of size $k$
(except possibly the last), in decreasing order of $|z_i|$.  For $j\ge1$,
\[
 \|z_{S_j}\|_2
 \le\frac1{\sqrt k}\|z_{S_{j-1}}\|_1,
\]
and hence, with $A=\|z\|_1/\sqrt k$,
\begin{equation}\label{eq:tailblocks}
 \sum_{j\ge1}\|z_{S_j}\|_2\le A,
 \qquad \|z_{S_0^c}\|_2\le A.
\end{equation}
Since $\Phi z=0$, RIP and the triangle inequality give
\[
 \sqrt{1-\delta}\,\|z_{S_0}\|_2
 \le\|\Phi z_{S_0}\|_2
 \le\sqrt{1+\delta}\sum_{j\ge1}\|z_{S_j}\|_2
 \le\sqrt{1+\delta}\,A.
\]
Using the orthogonality of the two coordinate pieces and
\eqref{eq:tailblocks},
\[
 \|z\|_2^2
 =\|z_{S_0}\|_2^2+\|z_{S_0^c}\|_2^2
 \le\left(\frac{1+\delta}{1-\delta}+1\right)A^2
 =\frac{2}{1-\delta}\frac{\|z\|_1^2}{k}.
\]
\end{proof}

\section{Improved asymptotic constant}

We use Theorem 8 of \cite{BLM}: if $U$ is Haar unitary, then
$\Phi=[I\ U]$ satisfies $(k,\delta)$-RIP with positive probability for all
sufficiently large $n$ provided $\delta<1$ and
\begin{equation}\label{eq:source-rip}
 n\ge\frac{256}{\delta^2}k\log\!\left(\frac{en}{k}\right).
\end{equation}

\begin{theorem}[Improved asymptotic lower bound]\label{thm:asymptotic}
For every sufficiently large $n$, there exists $U\in\mathrm U(n)$ such that
\[
 \ns(x)+\ns(Ux)\ge\frac3{20}\left\lfloor\frac n{5000}\right\rfloor
 \qquad(x\in\C^n\setminus\{0\}).
\]
Consequently,
\[
 \boxed{\displaystyle\liminf_{n\to\infty}c(n)\ge\frac3{100000}.}
\]
\end{theorem}

\begin{proof}
Take $\delta=7/10$ and $k=\lfloor n/5000\rfloor$.  The function
$\alpha\mapsto\alpha\log(e/\alpha)$ is increasing on $(0,1)$.  Therefore
\begin{align*}
 \frac{256}{\delta^2}\frac{k}{n}\log\!\left(\frac{en}{k}\right)
 &\le \frac{25600}{49}\frac1{5000}\log(5000e)\\
 &=\frac{256}{2450}\log(5000e)
 <\frac{256}{2450}(9.52)<1.
\end{align*}
Thus \eqref{eq:source-rip} holds.  The probability in \cite[Theorem 8]{BLM}
is positive once $n$ is sufficiently large, so there exists a unitary $U$
for which $\Phi=[I\ U]$ has this RIP.

For nonzero $x$, put $z=(Ux,-x)\in\ker[I\ U]$.  Lemma~\ref{lem:width}
gives
\[
 \frac{\|z\|_1^2}{\|z\|_2^2}\ge\frac{(1-\delta)k}{2}=\frac{3k}{20}.
\]
On the other hand, unitarity and the elementary inequality
$(a+b)^2/2\le a^2+b^2$ give
\begin{align*}
 \frac{\|z\|_1^2}{\|z\|_2^2}
 &=\frac{(\|x\|_1+\|Ux\|_1)^2}{2\|x\|_2^2}\\
 &\le\frac{\|x\|_1^2+\|Ux\|_1^2}{\|x\|_2^2}
 =\ns(x)+\ns(Ux).
\end{align*}
This proves the finite-$n$ assertion.  Divide by $n$ and let $n$ tend to
infinity.
\end{proof}

\section{Limitations and novelty check}

Theorem~\ref{thm:small} does not determine $c(4)$ or any later exact value.
Theorem~\ref{thm:asymptotic} improves only the lower constant; it does not
identify the asymptotically optimal one.  Numerical exploration of $F_4$,
$F_5$, and a small sample of Haar unitaries was used only to decide where
further exact arguments cease to be credible and supplies no proof claim.

A bounded literature search on 11 August 2026 used the exact source title and
arXiv id, the authors, the formula $\ns(x)+\ns(Ux)$, ``numerical sparsity
optimal unitary,'' and small-dimensional discrete Fourier $\ell_1$
uncertainty.  No later paper explicitly answering Problem 3 or recording the
values $c(2),c(3)$ was found.  Nearby finite-dimensional uncertainty papers use
different localization functionals.  Novelty confidence is moderate for the
exact small-dimensional values and moderate-to-high for the explicit
factor-$16.2$ improvement obtained from the source's own RIP theorem.

\paragraph{Human-review recommendation.}
Check the Fourier phase convention in \eqref{eq:q0}--\eqref{eq:q2}, the
largest-coordinate inequality in dimension three, the decreasing-block bound
\eqref{eq:tailblocks}, and the numerical inequality verifying
\eqref{eq:source-rip}.  If these pass, retain this as a substantial partial
answer to Problem 3.

\begin{thebibliography}{9}
\bibitem{BLM}
A.~S. Bandeira, M.~E. Lewis, and D.~G. Mixon,
\emph{Discrete uncertainty principles and sparse signal processing},
J. Fourier Anal. Appl. 24 (2018), 935--956; arXiv:1504.01014.
\end{thebibliography}

\end{document}
